%======================================================================
% CAPÍTULO 6: SANDRA CONCILIATION FILE (SCF)
%======================================================================
\chapter{Sandra Conciliation File (SCF)}

\section{Introducción al Módulo SCF}

\textbf{SCF (Sandra Conciliation File)} es el módulo de conciliación de archivos CSV integrado en Sentinel. Permite comparar dos archivos para identificar coincidencias y hallazgos.

\notacientifica{SCF utiliza un algoritmo de Hash Join en memoria, permitiendo comparar archivos de miles o millones de líneas eficientemente.}

\section{Uso del Sistema}

\subsection{Comando Conciliate}

El comando \texttt{conciliate} permite comparar archivos CSV directamente desde Sentinel:

\begin{verbatim}
sandra-sentinel conciliate \
  --comparison <archivo> \
  --comparison-columns <col1,col2,...> \
  --origin <archivo> \
  --origin-columns <col1,col2,...> \
  [--delimiter <char>] \
  [--output <dir>] \
  [--skip-header] \
  [--quiet]
\end{verbatim}

\notacientifica{La integración de SCF en Sentinel permite un flujo de trabajo unificado para procesamiento y conciliación.}

\subsection{Ejemplo Práctico}

\begin{caja}{Ejemplo de Uso}
\begin{verbatim}
# Comparar movimientos bancarios vs. sistema de nómina
sandra-sentinel conciliate \
  --comparison banco.csv \
  --comparison-columns 0,1,2 \
  --origin sistema.csv \
  --origin-columns 0,3,4 \
  --output resultados/
\end{verbatim}
\end{caja}

\section{Salida Generada}

\begin{caja}{Archivos de Resultado}
\begin{enumerate}
\item \textbf{coincidencias.txt}: Registros que coinciden entre ambos archivos.
\item \textbf{hallazgos.txt}: Registros del origen que no están en comparación.
\item \textbf{log.txt}: Errores de formato encontrados.
\item \textbf{reporte.txt}: Resumen con conteo de líneas.
\end{enumerate}
\end{caja}

\notacientifica{La separación de resultados facilita el flujo de trabajo de auditoría y conciliación.}

\section{Arquitectura de SCF}

SCF está diseñado con módulos independientes:

\begin{caja}{Módulos de SCF}
\begin{itemize}
\item \textbf{core/keys.rs}: Generación de claves de comparación.
\item \textbf{core/comparator.rs}: Motor de comparación Hash Join.
\item \textbf{io/mod.rs}: Lectura/escritura con buffering.
\item \textbf{report/mod.rs}: Generación de reportes.
\item \textbf{error.rs}: Manejo de errores personalizado.
\end{itemize}
\end{caja}

\notacientifica{La arquitectura modular permite reutilizar los componentes en diferentes contextos.}

\section{Algoritmo de Comparación}

SCF utiliza el algoritmo \textbf{Hash Join} para comparar archivos:

\begin{enumerate}
\item \textbf{Fase de Carga}: Lee el archivo de comparación y genera un HashSet de claves.
\item \textbf{Fase de Comparación}: Lee el archivo de origen y verifica cada registro contra el HashSet.
\item \textbf{Fase de Resultado}: Escribe los registros en archivos separados según el resultado.
\end{enumerate}

\notacientifica{La complejidad del algoritmo Hash Join es O(N+M), significativamente más eficiente que algoritmos de comparación línea a línea.}

\section{Integración con Sentinel}

SCF está integrado en Sentinel como módulo interno, permitiendo:

\begin{itemize}
\item Uso directo desde CLI de Sentinel.
\item Combinación con capacidades de cálculo del motor.
\item Flujo de trabajo unificado de procesamiento.
\item Integración con gRPC para procesamiento remoto.
\end{itemize}

\notacientifica{La integración elimina la necesidad de herramientas externas para la conciliación de archivos.}

\section{Conclusiones}

El módulo SCF complementa las capacidades de Sentinel proporcionando:

\begin{caja}{Resumen de Capacidades}
\begin{itemize}
\item \textbf{Comparación Eficiente}: Algoritmo Hash Join O(N+M).
\item \textbf{Formato Flexible}: Delimitador configurable.
\item \textbf{Claves Compuestas}: Múltiples columnas para identificación.
\item \textbf{Rendimiento Optimizado}: Buffered I/O, pre-alloc, paralelo.
\item \textbf{Salida Estructurada}: Archivos separados para auditoría.
\item \textbf{Integración Total}: Comando nativo en Sentinel.
\end{itemize}
\end{caja}

\notacientifica{SCF representa un ejemplo de ingeniería de sistemas especializada dentro del ecosistema Sandra.}
